\emph{सिद्धता-पद्धतींमुळे} वाक्यांमधील निष्पन्नतेचे आणि सुसंगतीचे लक्षणचित्रण करण्यासाठी निव्वळ विन्यासात्मक मार्ग मिळतात. \emph{नैसर्गिक निगमन} ही अशी एक सिद्धता-पद्धती आहे. त्यात \emph{निष्पत्ती} म्हणजे वाक्यांचा सांत वृक्ष असतो. त्याच्या पानांवरील वाक्ये \emph{गृहीतके} असतात किंवा बंद पद~$t$ साठी $\eq[t][t]$ या रूपाची वाक्ये असतात, जी $\Intro{\eq}$ नियमानुसार गृहीतकांशिवाय मिळतात. निष्पत्ती योग्य असण्यासाठी इतर प्रत्येक वाक्याचे अनेक \emph{निगमन नियमांपैकी} एका नियमानुसार योग्य रीतीने समर्थन केले पाहिजे. प्रत्येक आदिम संयोजक आणि प्रत्येक संख्यापकासाठी एक प्रवेशन नियम आणि एक विलोपन नियम असतो. उदाहरणार्थ, सूत्र~$!A$ चे समर्थन $\Elim{\lif}$ नियमानुसार केले असेल, तर त्याआधीची सूत्रे (म्हणजे \emph{आधारविधाने}) $!B \lif !A$ आणि $!B$ अशी असली पाहिजेत (एखाद्या~$!B$ साठी). काही निगमन नियम गृहीतके \emph{मुक्त करू} देतात. उदाहरणार्थ, $\Intro{\lif}$ वापरून $!B$ पासून $!A \lif !B$ काढले, तर आधारविधान~$!B$ पर्यंत नेणाऱ्या !!{derivation} मध्ये गृहीतक म्हणून आलेल्या~$!A$ च्या कोणत्याही आस्थिती मुक्त करता येतात. त्यांना एक खूण दिली जाते आणि ती त्या अनुमानापाशीही नोंदवली जाते.

अंतिम सूत्र~$!A$ असलेली अशी !!a{derivation} असेल की तिची सर्व गृहीतके मुक्त केलेली आहेत, तर $!A$ हे \emph{प्रमेय} आहे असे म्हणतो आणि~$\Proves !A$ असे लिहितो. अंतिम सूत्र~$!A$ असलेली अशी निष्पत्ती असेल की तिची मुक्त न केलेली सर्व गृहीतके संच~$\Gamma$ मध्ये आहेत, तर $!A$ हे~$\Gamma$ पासून \emph{सिद्धतायोग्य} आहे असे म्हणतो आणि $\Gamma \Proves !A$ असे लिहितो. $\Gamma \Proves \lfalse$ असेल, तर $\Gamma$ \emph{विसंगत} आहे असे म्हणतो, अन्यथा \emph{सुसंगत} आहे असे म्हणतो. या संकल्पना परस्परांशी संबंधित आहेत. उदाहरणार्थ, $\Gamma \Proves !A$ तेव्हा आणि केवळ तेव्हाच जेव्हा $\Gamma \cup \{\lnot !A\}$ विसंगत असतो. त्यांचा संबंधित चिन्हार्थविषयक संकल्पनांशीही संबंध आहे. उदाहरणार्थ, $\Gamma \Proves !A$ असेल, तर $\Gamma \Entails !A$. $\Gamma$ पासून जे काही सिद्ध करता येते ते~$\Gamma$ पासून तार्किकरीत्या निष्पन्न होते, या सिद्धता-पद्धतींच्या गुणधर्माला \emph{निर्दोषता} म्हणतात. \emph{निर्दोषता प्रमेय} निष्पत्तीतील अनुमानांच्या संख्येवर विगमन करून सिद्ध केले जाते. त्यात प्रत्येक अनुमानाबाबत असे दाखवले जाते की त्याचे प्रत्येक आधारविधान त्याच्यावर संपणाऱ्या उपनिष्पत्तीच्या मुक्त न केलेल्या गृहीतकांपासून तार्किकरीत्या निष्पन्न होत असेल, तर त्या अनुमानाचा निष्कर्ष संपूर्ण निष्पत्तीच्या मुक्त न केलेल्या गृहीतकांपासून तार्किकरीत्या निष्पन्न होतो.
